% !TeX root = Thesis.tex

% % % % % % % % % % Packages % % % % % % % % % %
\usepackage{cite}
\usepackage{appendix}
\usepackage{xcolor}
\usepackage{graphicx}
\usepackage{lmodern}
\usepackage{framed}
\usepackage[chapter]{algorithm}
\usepackage{verbatim}
\usepackage[english]{babel} %% Deutsche Bezeichnungen im Dokument
%\usepackage[ngerman,english]{babel} %% Englische Bezeichnungen im Dokument
\usepackage{hyperref}
\usepackage{amsmath}
\usepackage{amsfonts}
\usepackage{amssymb}
\usepackage{a4wide}
\usepackage[headsepline,footsepline]{scrlayer-scrpage}
\usepackage{microtype}
%\usepackage{fontspec} %% only if  using xelatex
\usepackage{float}
\usepackage[utf8]{inputenc}
\usepackage{stringstrings}
\usepackage{amsthm}
\usepackage[most]{tcolorbox}
\usepackage{prettyref}
\usepackage[margin=2.5cm, includefoot]{geometry}
\usepackage{wrapfig}
\usepackage[nolist]{acronym}

\begin{acronym}
    \acro{pet}[PET]{polyethylene terephthalate}
    \acro{pcb}[PCB]{printed circuit board}
\end{acronym}

\newtcbtheorem[auto counter,number within=chapter]{definitionbox}%
{Definition}{fonttitle=\bfseries\upshape, fontupper=\slshape,
    arc=0mm, colback=black!5!white,colframe=black!75!white}{definition}

\newtcbtheorem[auto counter,number within=chapter]{lemmabox}%
{Lemma}{fonttitle=\bfseries\upshape, fontupper=\slshape,
    arc=0mm, colback=blue!5!white,colframe=blue!75!black}{lemma}

\newtcbtheorem[auto counter,number within=chapter]{satzbox}%
{Satz}{fonttitle=\bfseries\upshape, fontupper=\slshape,
    arc=0mm, colback=red!5!white,colframe=red!75!black}{satzbox}

\newrefformat{lemma}{Lemma~\ref{#1}}
\newrefformat{sec}{Section~\ref{#1}}
\newrefformat{subsec}{Section~\ref{#1}}
\newrefformat{definition}{Definition~\ref{#1}}
\newrefformat{fig}{Figure~\ref{#1}}
\newrefformat{cap}{Chapter~\ref{#1}}

% % % % % % % % % % oft verwendete Farben % % % % % % % % % %
\definecolor{black}{rgb}{0,0,0}
\definecolor{darkgrey}{rgb}{.4,.4,.4}
\definecolor{grey}{rgb}{.7,.7,.7}
\definecolor{white}{rgb}{1,1,1}

\definecolor{blue}{rgb}{0,0,1}
\definecolor{red}{rgb}{1,0,0}
\definecolor{green}{rgb}{0,1,0}

\definecolor{midblue}{rgb}{0,0,.7}
\definecolor{midred}{rgb}{.7,0,0}
\definecolor{midgreen}{rgb}{0,.7,0}

\definecolor{darkblue}{rgb}{0,0,.4}
\definecolor{darkred}{rgb}{.4,0,0}
\definecolor{darkgreen}{rgb}{0,.4,0}

\definecolor{orange}{rgb}{1,.4,0}
\definecolor{cyan}{rgb}{0,1,.8}
\definecolor{purple}{rgb}{.6,0,.6}
\definecolor{lavender}{rgb}{.6,.6,1}
\definecolor{yellow}{rgb}{1,.8,.2}



% % % % % % % % % %  RWTH Colors % % % % % % % % % %  
\definecolor{rwthgrey75}{RGB}{100,101,103}
\definecolor{rwthgrey50}{RGB}{156,158,159}
\definecolor{rwthgrey25}{RGB}{207,209,210}
\definecolor{rwthgrey10}{RGB}{236,237,237}

\definecolor{rwthblue100}{HTML}{00549F}
\definecolor{rwthblue75}{HTML}{407FB7}
\definecolor{rwthblue50}{HTML}{8EBAE5}
\definecolor{rwthblue25}{HTML}{C7DDF2}
\definecolor{rwthblue10}{HTML}{E8F1FA}

% % % % % % % % % %  FH, IT Center and MATSE Colors % % % % % % % % % % 
\definecolor{fhgreen}{HTML}{00B1AC}
\definecolor{itcorange}{HTML}{F6A800}
\definecolor{matseblue}{HTML}{00549F}


% % % % % % % % % %  Farboptionen % % % % % % % % % % 
%% Screen ist bunt, Print ist schwarz. Passendes einkommentieren.

% Screen
\hypersetup{
  pdftitle={Report},
  pdfauthor={Julius Gummersbach, 114015019},
  pdfsubject={An Origami Continuum Robot Capable of Precise Motion Through Torsionally Stiff Body and Smooth Inverse Kinematics},
  colorlinks=true,
  linkcolor=midblue,
  citecolor=midgreen,
  urlcolor=midred
}

% Print
%\hypersetup{
%  pdftitle={Seminararbeit},
%  pdfauthor={Name, Matrikelnummer},
%  pdfsubject={Titel der Arbeit},
%  colorlinks=true,
%  linkcolor=black,
%  citecolor=black,
%  urlcolor=black
%}

% % % % % % % % % %  Kopf- und Fußzeile % % % % % % % % % % 
\pagestyle{plain}
\automark[section]{chapter}
\cfoot{\vspace{3pt}\phantom{ }}


% % % % % % % % % %  Umgebung für Beispiele % % % % % % % % % % 
\newlength{\leftbarwidth}
\setlength{\leftbarwidth}{1pt}
\newlength{\leftbarsep}
\setlength{\leftbarsep}{10pt}
\newlength{\leftbarheight}
\setlength{\leftbarheight}{10pt}

\newcommand*{\leftbarcolorcmd}[1]{\color{#1}}% as a command to be more flexible


%%% Useful for marking examples, which are inserted into the text. Usage see below.
\newenvironment{example}[1][black]{%
    \def\FrameCommand{{\leftbarcolorcmd{#1}{\hspace{3pt}  \vrule width \leftbarwidth \relax\hspace {\leftbarsep}}}}%
    \MakeFramed {\vspace{2pt} \advance \hsize -\width \FrameRestore }%
\vspace{-0.5em}}{ %
    \vspace{2pt} \endMakeFramed
}

\renewenvironment{leftbar}[1][black]{%
    \def\FrameCommand{{\leftbarcolorcmd{#1}{\hspace{3pt}  \vrule width \leftbarwidth \relax\hspace {\leftbarsep}}}}%
    \MakeFramed {\vspace{2pt} \advance \hsize -\width \FrameRestore }%
\vspace{-1.5em}}{ %
    \vspace{2pt} \endMakeFramed
}

% % % % % % % % % %  Generelle Einstellungen im Dokument % % % % % % % % % % 
\setlength{\parindent}{0pt}


% % % % % % % % % %  Eigene Kommandos % % % % % % % % % %
\newcommand\forceblankpage{\newpage\thispagestyle{empty}\phantom{ }\newpage}

% % % % % % % % % %  Einstellungen KOMA-Script % % % % % % % % % % 
\KOMAoptions{cleardoublepage=empty, parskip=half}